\chapter{Sống Cùng Màn Hình Mọi Nơi}
\markboth{Sống Cùng Màn Hình Mọi Nơi}{Sống Cùng Màn Hình Mọi Nơi}

\begin{center}
	\textit{\color{bookprimary}``Nếu đi vệ sinh cũng phải xem, đi bộ cũng phải xem, tập thể dục cũng phải xem, thì điện thoại không còn là công cụ nữa. Nó đã thành dây xích.''}

	\textit{\color{bookprimary}``If you must watch while using the toilet, walking, and exercising, then the phone is no longer a tool. It has become a chain.''}
\end{center}

\ngancach

\section{Ngay Cả Lúc Cơ Thể Đang Cần Thở}
\begin{truyen}{Ngay Cả Lúc Cơ Thể Đang Cần Thở}{Even When the Body Needs Air}
\chuhoa{C}{ó bạn vừa đi bộ thể dục quanh sân trường vừa dán mắt vào điện thoại.}
\textit{(A student walks around the school yard for exercise while staring at the phone.)}

Thầy hỏi: ``Em đang tập cái gì?''
\textit{(The teacher asks: ``What exactly are you training?'')}

Bạn ấy cười: ``Dạ em vừa tập vừa thư giãn mà thầy.''
\textit{(The student laughs: ``I am exercising and relaxing at the same time.'')}

\teacherpair{Không. Em đang làm cả hai thứ cùng dở đi. Tập thể dục là lúc máu phải chạy, hơi thở phải sâu, đầu óc phải có cơ hội rời bớt kích thích. Còn em lại kéo thêm một cái màn hình vào. Cơ thể em đang vận động, nhưng đầu óc em vẫn bị buộc bằng một sợi dây rất ngắn.}{``No. You are doing both things more poorly. Exercise is when blood should move, breathing should deepen, and the mind should get a chance to step away from stimulation. But you drag a screen right into it. Your body is moving, while your mind remains tied by a very short leash.''}

Thầy nói tiếp: ``Nhiều người bây giờ không hề được nghỉ thật. Họ chỉ thay đổi tư thế trong khi vẫn tiếp tục tiêu thụ kích thích. Thế nên cơ thể có thể mệt hơn mà đầu óc vẫn không hề nhẹ hơn.''
\textit{(The teacher continues: ``Many people today never truly rest. They only change posture while continuing to consume stimulation. That is why the body may become more tired while the mind does not feel lighter at all.'')}
\end{truyen}

\begin{baihoc}
Có những hoạt động đáng lẽ phải cứu đầu óc khỏi màn hình.
\textit{(Some activities are supposed to rescue the mind from the screen.)}

Nếu em cũng nhét màn hình vào nốt những lúc đó, em đang tự cắt mất đường hồi sức của chính mình.
\textit{(If you also stuff the screen into those moments, you are cutting off your own path to recovery.)}
\end{baihoc}

\begin{tuhocnhanh}
1. Em có đang mang màn hình vào cả những lúc cơ thể cần được tự do không? \textit{(Are you bringing the screen even into moments when the body needs freedom?)}

2. Hoạt động nào trong ngày của em đáng lẽ nên không có điện thoại nhất? \textit{(Which daily activity of yours most deserves to be phone-free?)}

3. Em có còn nhớ cảm giác đi bộ mà đầu óc không bị kéo đi đâu nữa không? \textit{(Do you still remember what it feels like to walk without your mind being pulled away?)}
\end{tuhocnhanh}

\begin{tongketchuong}
Đem màn hình vào mọi nơi không phải là sống hiện đại hơn.
\textit{(Dragging the screen into every place is not a more modern life.)}

Nhiều khi đó chỉ là một kiểu lệ thuộc được bọc trong vẻ tiện lợi.
\textit{(Often it is only dependence wrapped in convenience.)}
\end{tongketchuong}

\ngancach